\documentclass[11pt]{article}
\usepackage[margin=1in]{geometry}
\usepackage{amsmath,amssymb}
\usepackage{booktabs}
\usepackage{hyperref}
\usepackage{xcolor}
\usepackage{listings}

\lstset{basicstyle=\ttfamily\footnotesize, frame=single, breaklines=true,
        columns=fullflexible, xleftmargin=1em}

% Long \texttt{} identifiers (file paths, test names) cannot hyphenate;
% give TeX room to stretch rather than overflowing the margin.
\emergencystretch=4em

\newcommand{\pn}{\phi_n}
\newcommand{\dd}{\,\mathrm{d}}
\newcommand{\OK}{{\color{blue}\textbf{matches}}}
\newcommand{\BAD}{{\color{red}\textbf{HALF}}}

\title{Audit of the repulsion coefficient in the CUDA phase-field cell simulation}
\author{Code review note}
\date{2026-07-28}

\begin{document}
\maketitle

\begin{abstract}
\noindent
The interaction (repulsion) coefficient in the $\phi$ evolution equation is
implemented at $30\kappa/\lambda^2$, whereas Palmieri \emph{et al.}\ (2015)
Eq.~(S15) requires $60\kappa/\lambda^2$ --- a factor of $2$ too small. The
velocity coefficient $60\kappa/(\xi\lambda^2)$ is correct. Every other term in
the right-hand side matches the paper exactly. The missing factor is the
double-counting of the ordered pair sum $\sum_{n,m\neq n}$ in the interaction
free energy, not the mobility. The discrepancy has been confirmed by direct
black-box measurement of the compiled CUDA binary
(Sec.~\ref{sec:measurement}). Net effect: the simulation reproduces Palmieri's
model at $\kappa_{\rm eff}=5$, $\xi_{\rm eff}=750$ rather than the reported
$\kappa=10$, $\xi=1500$.
\end{abstract}

\section{Provenance and confidence}

The model equations below were read from the local PDFs
\texttt{papers/10.1038\_srep11745.pdf} (main text) and
\texttt{papers/srep11745\_SI.pdf} (supplementary), extracted with
\texttt{pypdf} in both default and \texttt{layout} mode. PDF mathematical
typesetting does not extract cleanly, so the confidence in each reading is
stated explicitly:

\begin{itemize}
\item \textbf{Eq.~(S15) --- high confidence.} The \texttt{layout}-mode
  extraction unambiguously shows a single prefactor $-30/\tilde\lambda^2$
  applied to a bracket containing \emph{both} the double-well term and the
  term $2\tilde\kappa\sum_{m\neq n}\phi_n\phi_m^2$. This is the operative
  equation and it is the basis of the whole finding.
\item \textbf{Eq.~(10) --- high confidence.} Prefactor $30\kappa/\lambda^2$ and
  the ordered sum $\sum_{n,m\neq n}$ both extract cleanly.
\item \textbf{Eq.~(12) --- high confidence.} Coefficient $60\kappa/(\xi\lambda^2)$
  extracts cleanly.
\item \textbf{Eq.~(7) volume term --- medium confidence.} The extraction is
  garbled, but shows two $\pi$ and two $R$ tokens, consistent with the
  normalised form $\frac{\mu_n}{\pi R^2}(\pi R^2 - \int\phi_n^2)^2$. This is
  independently confirmed by SI Table~I, where $\mu = \gamma_0\tilde\mu /
  (\pi\tilde R^2\Delta^4)$, and by Eq.~(S15), whose volume term is
  $-\frac{2\tilde\mu}{\pi\tilde R^2}$.
\end{itemize}

\noindent
Crucially, the finding does \emph{not} depend on the medium-confidence
reading. Sec.~\ref{sec:crosschecks} gives two independent cross-checks, and
Sec.~\ref{sec:measurement} gives a direct numerical measurement of the
compiled binary.

\section{Palmieri's model as stated}

All quantities are dimensionless; the paper states that Eq.~(1) onward is
written in dimensionless units.

\subsection{Free energy}

Single-cell part, Eq.~(7):
\begin{equation}
\mathcal{F}_0 = \sum_n\left\{
  \int\!\!\!\int \dd x\dd y \left[
     \gamma_n(\nabla\pn)^2
   + \frac{30\gamma_n}{\lambda^2}\,\pn^2(1-\pn)^2
  \right]
  + \frac{\mu_n}{\pi R^2}\left(\pi R^2 - \int\!\!\!\int \dd x\dd y\,\pn^2\right)^{\!2}
\right\}
\label{eq:F0}
\end{equation}

Total free energy, Eq.~(9), and the interaction part, Eq.~(10):
\begin{equation}
\mathcal{F} = \mathcal{F}_0 + \mathcal{F}_{\rm int},
\qquad
\mathcal{F}_{\rm int} = \int\!\!\!\int \dd x\dd y
   \sum_{n,\,m\neq n} \frac{30\kappa}{\lambda^2}\,\pn^2\phi_m^2 .
\label{eq:Fint}
\end{equation}

\begin{quote}
\textbf{Note the sum in Eq.~\eqref{eq:Fint} is \emph{ordered}}: it runs over
all pairs $(n,m)$ with $m\neq n$, so every unordered pair is counted twice.
This is the origin of the factor under dispute.
\end{quote}

\subsection{Dynamics}

Equation of motion, Eq.~(1), with mobility $M=\tfrac12$:
\begin{equation}
\frac{\partial\pn}{\partial t} + \mathbf{v}_n\cdot\nabla\pn
  = -\frac{1}{2}\,\frac{\delta\mathcal{F}}{\delta\pn}.
\label{eq:eom}
\end{equation}

Velocity decomposition, Eq.~(11), and the interaction velocity, Eq.~(12):
\begin{equation}
\mathbf{v}_n = \mathbf{v}_{n,I} + \mathbf{v}_{n,A},
\qquad
\mathbf{v}_{n,I} = \frac{60\kappa}{\xi\lambda^2}
  \int\!\!\!\int \dd x\dd y\;\pn\,(\nabla\pn)\!\!\sum_{m\neq n}\!\phi_m^2 .
\label{eq:vI}
\end{equation}

Run-and-tumble reorientation, Eq.~(13):
\begin{equation}
P(t_r)\dd t_r = \frac{1}{\tau}e^{-t_r/\tau}\dd t_r .
\label{eq:rtp}
\end{equation}

\subsection{The PDE that is actually integrated}

SI Eq.~(S15) is Eq.~\eqref{eq:eom} with the functional derivative carried out.
\textbf{This is the single most important equation in this document:}
\begin{equation}
\boxed{
\begin{aligned}
\frac{\partial\pn}{\partial t} + \mathbf{v}_n\cdot\nabla\pn
&= \gamma_n\nabla^2\pn
 - \frac{30}{\lambda^2}\left[
      \gamma_n\pn(1-\pn)(1-2\pn)
    + 2\kappa\!\!\sum_{m\neq n}\!\pn\phi_m^2
  \right]\\[2pt]
&\quad - \frac{2\mu}{\pi R^2}\,\pn\left[
      \int\!\!\!\int \dd x\dd y\,\pn^2 - \pi R^2
  \right]
\end{aligned}}
\label{eq:S15}
\end{equation}
Both the double-well term and the repulsion term sit inside \emph{one} bracket
sharing \emph{one} prefactor $-30/\lambda^2$. Expanding:
\begin{align}
\text{double-well contribution} &= -\frac{30\gamma_n}{\lambda^2}\pn(1-\pn)(1-2\pn),\\
\text{repulsion contribution}   &= -\frac{30}{\lambda^2}\cdot 2\kappa \cdot \pn\!\!\sum_{m\neq n}\!\phi_m^2
                                 = -\frac{60\kappa}{\lambda^2}\,\pn\!\!\sum_{m\neq n}\!\phi_m^2 .
\label{eq:target}
\end{align}

\subsection{Parameters (Table 1)}

\begin{center}
\begin{tabular}{lcccccccc}
\toprule
& $\gamma_n$ & $\lambda$ & $R$ & $\kappa$ & $\mu$ & $\tau$ & $v_A$ & $\xi$ \\
\midrule
value & 0.35 / 1 & 7 & 49 & 10 & 1 & $10^4$ & $10^{-2}$ & $1.5\times10^3$ \\
cell type & cancer / normal & \multicolumn{7}{c}{all} \\
\bottomrule
\end{tabular}
\end{center}

\section{Verification that Eq.~(S15) follows from Eqs.~\eqref{eq:F0}--\eqref{eq:eom}}

This confirms the paper is internally self-consistent, and isolates where the
factor of $2$ comes from.

\paragraph{1. Gradient term.}
$\dfrac{\delta}{\delta\pn}\displaystyle\int\gamma_n(\nabla\pn)^2 = -2\gamma_n\nabla^2\pn$.
Times $-\tfrac12$: $\;+\gamma_n\nabla^2\pn$. \OK\ Eq.~\eqref{eq:S15}.

\paragraph{2. Double well.}
Using $\dfrac{d}{d\phi}\left[\phi^2(1-\phi)^2\right] = 2\phi(1-\phi)(1-2\phi)$,
\[
\frac{\delta}{\delta\pn}\int\frac{30\gamma_n}{\lambda^2}\pn^2(1-\pn)^2
 = \frac{60\gamma_n}{\lambda^2}\pn(1-\pn)(1-2\pn).
\]
Times $-\tfrac12$: $\;-\dfrac{30\gamma_n}{\lambda^2}\pn(1-\pn)(1-2\pn)$. \OK

\paragraph{3. Volume constraint.}
With $V_n=\int\pn^2$ and $A_0=\pi R^2$,
\[
\frac{\delta}{\delta\pn}\;\frac{\mu}{A_0}(A_0-V_n)^2
 = \frac{\mu}{A_0}\cdot 2(A_0-V_n)\cdot(-2\pn)
 = -\frac{4\mu}{A_0}(A_0-V_n)\pn .
\]
Times $-\tfrac12$: $\;+\dfrac{2\mu}{A_0}(A_0-V_n)\pn = -\dfrac{2\mu}{\pi R^2}\pn\!\left[V_n-\pi R^2\right]$. \OK

\paragraph{4. Repulsion --- where the factor of 2 lives.}
In the ordered sum of Eq.~\eqref{eq:Fint}, index $n$ appears in \emph{both}
slots. Collecting all $\pn$-dependent terms:
\[
\mathcal{F}_{\rm int}\Big|_{\pn}
 = \frac{30\kappa}{\lambda^2}\Big[
     \underbrace{\sum_{m\neq n}\pn^2\phi_m^2}_{n\ \text{in first slot}}
   + \underbrace{\sum_{m\neq n}\phi_m^2\pn^2}_{n\ \text{in second slot}}
   \Big]
 = 2\,\frac{30\kappa}{\lambda^2}\,\pn^2\!\!\sum_{m\neq n}\!\phi_m^2 .
\]
Hence
\[
\frac{\delta\mathcal{F}_{\rm int}}{\delta\pn}
 = 2\cdot\frac{30\kappa}{\lambda^2}\cdot 2\pn\!\!\sum_{m\neq n}\!\phi_m^2
 = \frac{120\kappa}{\lambda^2}\,\pn\!\!\sum_{m\neq n}\!\phi_m^2 ,
\]
and times $-\tfrac12$:
\begin{equation}
\boxed{\;-\frac{60\kappa}{\lambda^2}\,\pn\!\!\sum_{m\neq n}\!\phi_m^2\;}
\end{equation}
which is exactly the $2\tilde\kappa$ term printed inside the S15 bracket,
Eq.~\eqref{eq:target}. \OK

\paragraph{5. Consistency of the velocity.}
Only $\mathcal{F}_{\rm int}$ survives in $\int(\delta\mathcal{F}/\delta\pn)\nabla\pn$
(the single-cell terms are total derivatives and integrate to zero). With the
same mobility $M=\tfrac12$,
\[
\mathbf{v}_{n,I} = \frac{1}{2\xi}\int\frac{\delta\mathcal{F}}{\delta\pn}\nabla\pn
 = \frac{1}{2\xi}\cdot\frac{120\kappa}{\lambda^2}\int\pn(\nabla\pn)\!\!\sum_{m\neq n}\!\phi_m^2
 = \frac{60\kappa}{\xi\lambda^2}\int\pn(\nabla\pn)\!\!\sum_{m\neq n}\!\phi_m^2 ,
\]
which is Eq.~\eqref{eq:vI} exactly. \OK\ \emph{The paper is internally
consistent: Eq.~(S15) and Eq.~(12) both follow from the same
$M\,\delta\mathcal{F}/\delta\phi$.}

\section{What the code implements}

\subsection{CUDA}

Coefficient definitions, \texttt{include/types.cuh:26--36}:
\begin{lstlisting}
template <typename T> __host__ __device__ inline T bulk_coeff(T lambda) {
    return T(30) / (lambda * lambda); }
template <typename T> __host__ __device__ inline T interaction_coeff(T kappa, T lambda) {
    return T(30) * kappa / (lambda * lambda); }
template <typename T> __host__ __device__ inline T motility_coeff(T kappa, T xi, T lambda) {
    return T(60) * kappa / (xi * lambda * lambda); }
\end{lstlisting}

Use, \texttt{src/kernels.cu:505--506}, \texttt{546}, \texttt{548}:
\begin{lstlisting}
505:  dwC_s  = gam * bulk_coeff(lambda_);              // 30*gamma/lambda^2
506:  repC_s = interaction_coeff(kappa, lambda_);      // 30*kappa /lambda^2
...
546:  float rhs = gam*lap - dwC*dw + volC*c - repC*c*Soth;
547:  float adv = vx*gx + vy*gy;
548:  outp[idx] = c + dt * (rhs - adv);
\end{lstlisting}
There is no further scaling between line 506 and line 548: \texttt{rhs} is
$\partial\phi/\partial t$ directly.

\subsection{Reference implementations}

Both CPU references use the same values, so neither is an independent oracle.
From \texttt{tests/python/} \texttt{cpu\_reference.py}:
\begin{lstlisting}
132:  two_keff = 60.0 * p.kappa / (p.lambd ** 2)        # 12.245
165:  repulsion  = two_keff * phi * S
166:  var_deriv  = -tg * lap + bulk + constraint + repulsion
170:  phi_new    = phi + p.dt * (-0.5 * var_deriv - advection)   # -> 6.122
\end{lstlisting}
\texttt{rust/cpu\_ref/src/sim.rs}: lines 199, 342, 343, 345 are structurally
identical. Note the trap: line 132 in isolation \emph{looks} correct, because
$60\kappa/\lambda^2$ is precisely the S15 coefficient --- but it is placed in
\texttt{var\_deriv}, which line 170 multiplies by $-0.5$.

\subsection{Side-by-side}

\begin{center}\small
\begin{tabular}{lccc}
\toprule
Term in $\partial\phi/\partial t$ & Palmieri (S15) & Code & Status \\
\midrule
Laplacian     & $+\gamma$                  & $+\gamma$                  & \OK \\
Double well   & $-30\gamma/\lambda^2 = 0.612245$ & $0.612245$ & \OK \\
Volume        & $-2\mu/\pi R^2$            & $-2\mu/\pi R^2$            & \OK \\
Repulsion     & $-60\kappa/\lambda^2 = 12.244898$ & $6.122449$ & \BAD \\
Velocity      & $60\kappa/\xi\lambda^2 = 8.16327\times10^{-3}$ & same & \OK \\
\bottomrule
\end{tabular}
\end{center}
(numerical values at $\gamma=1,\ \kappa=10,\ \lambda=7,\ \mu=1,\ \xi=1500$.)

\section{Two cross-checks independent of the free-energy convention}
\label{sec:crosschecks}

Both of the following hold regardless of how one normalises
$\mathcal{F}_{\rm int}$, so they cannot be evaded by a convention argument.

\paragraph{Check A: bulk/repulsion ratio.}
In Eq.~\eqref{eq:S15} the double-well and repulsion terms share the single
prefactor $-30/\lambda^2$. Their ratio is therefore fixed by the paper:
\[
\frac{\text{repulsion coeff}}{\text{bulk coeff}} = \frac{2\kappa}{\gamma} = 20
\quad\text{(at }\kappa=10,\gamma=1\text{)},
\qquad\text{code: } \frac{30\kappa/\lambda^2}{30\gamma/\lambda^2} = \frac{\kappa}{\gamma} = 10 .
\]
In the code these two coefficients are two terms of the \emph{same sum on the
same line} (\texttt{kernels.cu:546}), multiplied by the same \texttt{dt} on
line 548. Any global factor --- the mobility, a units convention --- scales
both equally and cancels in the ratio. Since the bulk coefficient
independently matches S15, the repulsion coefficient is short by $2$.

\paragraph{Check B: repulsion/velocity ratio.}
The $\phi$-equation repulsion coefficient and the velocity coefficient are the
same object $M\,\delta\mathcal{F}_{\rm int}/\delta\phi$, one of them divided by
$\xi$. They must therefore differ by exactly $\xi$:
\[
\text{Palmieri: } \frac{60\kappa/\lambda^2}{60\kappa/\xi\lambda^2} = \xi = 1500 ,
\qquad
\text{code: } \frac{30\kappa/\lambda^2}{60\kappa/\xi\lambda^2} = \frac{\xi}{2} = 750 .
\]

Both checks independently demand $\text{repC} = 60\kappa/\lambda^2 = 12.2449$.

\section{Direct measurement of the compiled binary}
\label{sec:measurement}

To remove any reliance on reading source code, the coefficient was measured
from the binary's own output. Sixteen cells at $\rho=0.9$ were equilibrated to
firm contact, then resumed with $\xi=10^{12}$ so that the interaction velocity
(hence all advection) is $\sim10^{-11}$ and the dynamics is pure gradient flow.
Two checkpoints one $\mathrm{d}t$ apart give $\partial\phi/\partial t$ by finite
difference. Every term except the repulsion is independently verified, so one
solves pointwise:
\begin{equation}
C = -\,\frac{\left.\dfrac{\partial\phi}{\partial t}\right|_{\rm measured}
   - \left[\gamma\nabla^2\phi - \dfrac{30\gamma}{\lambda^2}\phi(1-\phi)(1-2\phi)
     + \dfrac{2\mu}{A_0}(A_0-V)\phi\right]}
  {\phi\,S_{\rm other}} .
\end{equation}

\begin{lstlisting}
  cell  0:   1494 px   C =   6.1224  (IQR 6.1224..6.1225)
  cell  1:   1738 px   C =   6.1224  (IQR 6.1224..6.1225)
  cell  2:   1379 px   C =   6.1224  (IQR 6.1224..6.1225)
  cell  3:   1176 px   C =   6.1224  (IQR 6.1224..6.1225)

  pooled over 16 cells, 23265 pixels
  MEASURED repulsion coefficient C = 6.12245
    candidate 30k/l^2  =  6.12245   ratio C/cand = 1.00000
    candidate 60k/l^2  = 12.24490   ratio C/cand = 0.50000
\end{lstlisting}

\noindent
Agreement with $30\kappa/\lambda^2$ to six significant figures over
$2.3\times10^4$ pixels. \textbf{The binary evaluates half of the
Palmieri coefficient.}

\section{Consequences and remedy}

\subsection{Equivalent parameter set}

Since $\kappa$ and $\xi$ enter \emph{only} through \texttt{interaction\_coeff}
and \texttt{motility\_coeff} (verified by exhaustive grep), the implemented
model is exactly Palmieri's model evaluated at
\begin{equation}
\kappa_{\rm eff} = 5, \qquad \xi_{\rm eff} = 750,
\end{equation}
while the runs are labelled $\kappa=10$, $\xi=1500$. Check:
$60\cdot5/\lambda^2 = 6.1224$ \checkmark\ and
$60\cdot5/(750\lambda^2) = 8.163\times10^{-3}$ \checkmark.

\subsection{Zero-code-change workaround}

By the same token, running the \emph{current unmodified binary} with
\begin{center}
\texttt{--kappa 20 --xi 3000}
\end{center}
gives repulsion $30\cdot20/\lambda^2 = 12.2449$ \checkmark\ and motility
$60\cdot20/(3000\lambda^2) = 8.163\times10^{-3}$ \checkmark\ --- i.e.\ exactly
Palmieri's intended physics. This permits a paired A/B comparison quantifying
the impact on observables before any code is touched, and serves as a
validation target for any patch.

\subsection{Code fix}

\begin{center}\footnotesize
\begin{tabular}{lll}
\toprule
File & Current & Fixed \\
\midrule
\texttt{src/kernels.cu:506}  & \texttt{interaction\_coeff(kappa, lambda\_)} & \texttt{2.0f * interaction\_coeff(...)} \\
\texttt{cpu\_reference.py:132} & \texttt{60.0 * kappa / lambd**2} & \texttt{120.0 * kappa / lambd**2} \\
\texttt{rust/cpu\_ref/sim.rs:199} & \texttt{60.0 * kappa / lambd**2} & \texttt{120.0 * kappa / lambd**2} \\
\bottomrule
\end{tabular}
\end{center}
\texttt{motility\_coeff} must \textbf{not} be changed.

\subsection{Why no existing test caught this}

The two tests that act as physics oracles ---
\texttt{test\_isolated\_cell\_energy\_decreases} and
\texttt{test\_single\_cell\_pde\_residual\_small} --- are both
\textbf{single-cell}. The latter states in a comment ``No interaction term
(single cell)''. The repulsion coefficient is the only coefficient that
requires $\geq 2$ cells to exercise, so no oracle test ever evaluates it. All
other physics tests compare the CUDA code against
\texttt{cpu\_reference.py} or the Rust reference, which encode the same
derived constants --- circular by construction.

Recommended additions: (i) a convention-free assertion that
$\text{repC}/\text{motility\_coeff} = \xi$; and (ii) a two-cell variational
test that writes $\mathcal{F}$ explicitly and differentiates it
\emph{numerically}, asserting $\partial_t\phi = -\tfrac12\,\delta
\mathcal{F}/\delta\phi$ in the overlap region.

\subsection{Impact on motility results}

\begin{itemize}
\item \textbf{Dilute / isolated cells: unaffected.} $\mathbf v_{n,I}\to0$, so
  $\mathbf v = v_A\hat{\mathbf p}$ and Palmieri Eq.~(14),
  $\langle\Delta x^2\rangle = 2v_A^2\tau^2[t/\tau - 1 + e^{-t/\tau}]$ with
  $D_{\rm eff}=v_A^2\tau/2$, holds exactly.
\item \textbf{Confluent monolayer: affected.} Velocities are correct for a
  given configuration, but the configurations visited are not: cells
  interpenetrate roughly twice as far before generating the same restoring
  force. Palmieri chose $\kappa$ large specifically so cells ``deform when
  brought together rather than overlap''; halving it moves the system toward
  the regime he tuned against.
\item Deformation observables ($L_n$, and hence $p_{\rm eff}=2\sqrt{\pi}L$)
  are compressed toward unity; the soft/normal contrast shrinks in magnitude.
\item Caging is weaker, so $D_{\rm eff}$ in the monolayer runs high and the
  jamming/glass transition $\rho_c$ shifts upward --- the principal risk for
  any finite-size-scaling analysis.
\item Speed bursts, whose mechanism is neighbour pinching during T1 events
  with force $\propto\kappa$, are weakened.
\item Existing equilibration checkpoints relaxed into a more-overlapped,
  less-deformed packing, so they are not valid initial conditions for a
  corrected binary.
\item Qualitative comparative conclusions driven by the elasticity contrast
  $\gamma_c/\gamma_n=0.35$ --- which is implemented correctly --- are likely
  to survive; quantitative comparison against Palmieri's published numbers is
  not.
\end{itemize}

\section{Manuscript equations}

For completeness, the free energy as typeset in
\texttt{nibi\_sync/manuscript.tex} does not reproduce
Eqs.~\eqref{eq:F0}--\eqref{eq:Fint} either, in two of three terms:

\begin{center}\small
\begin{tabular}{lll}
\toprule
& Manuscript (as typeset) & vs.\ Palmieri \\
\midrule
$F_{\rm int}$ (line 278) & $\gamma_i|\nabla\phi_i|^2 + \tfrac{30\gamma_i}{\lambda^2}\phi_i^2(1-\phi_i)^2$ & identical \OK \\
$F_{\rm vol}$ (line 289) & $\mu\sum_i\left(\int\phi_i^2 - A_0\right)^2$ & missing $1/A_0$ \\
$F_{\rm rep}$ (line 298) & $\kappa\sum_{i<j}\int\phi_i^2\phi_j^2$ & missing $60/\lambda^2$ \\
\bottomrule
\end{tabular}
\end{center}

\noindent
Since $\sum_{n,m\neq n} = 2\sum_{i<j}$, Palmieri's Eq.~(10) equals
$\frac{60\kappa}{\lambda^2}\sum_{i<j}$; matching it with the manuscript's form
would require $\kappa = 12.24$, whereas the manuscript's parameter table
(line 369) quotes Palmieri's $\kappa=10$. The manuscript states the model is
used ``without modification'', so these two equations should be corrected to
\[
F_{\rm rep} = \frac{30\kappa}{\lambda^2}\sum_{i,\,j\neq i}\int\phi_i^2\phi_j^2 ,
\qquad
F_{\rm vol} = \frac{\mu}{A_0}\sum_i\left(\int\phi_i^2 - A_0\right)^2 .
\]
Note that the manuscript's error and the code's error are in \emph{different}
directions ($-10\phi S$ implied by the manuscript, $-6.12$ evaluated by the
code, $-12.24$ required by Palmieri), so they are independent fixes.

\end{document}
